% Generated draft: requires mathematical and source review.
\documentclass[11pt]{article}
\usepackage[margin=25mm]{geometry}
\usepackage{fontspec}
\setmainfont{Latin Modern Roman}
\usepackage{amsmath,amssymb,mathtools,unicode-math}
\setmathfont{Latin Modern Math}
\usepackage{enumitem,needspace,xurl,hyperref}
\setlength{\parindent}{0pt}
\setlength{\parskip}{5pt}
\setlength{\emergencystretch}{3em}
\newcommand{\N}{\mathbb{N}}
\newcommand{\Z}{\mathbb{Z}}
\newcommand{\Q}{\mathbb{Q}}
\newcommand{\R}{\mathbb{R}}
\newcommand{\C}{\mathbb{C}}
\newcommand{\F}{\mathbb{F}}
\newcommand{\A}{\mathbb{A}}
\newcommand{\PP}{\mathbb{P}}
\newcommand{\E}{\mathbb{E}}
\newcommand{\eps}{\varepsilon}
\newcommand{\dd}{\,\mathrm{d}}
\begin{document}
\section*{UnsolvedMath ID: 30002455}
\textbf{Source problem code:} catalogue entry
\subsection{Definitions and mathematical statement}
Let \(\mathbb P\) be a coefficient semifield: an abelian multiplicative group equipped with a commutative, associative auxiliary addition \(\oplus\) that distributes over multiplication. Write \(\mathbb Z\mathbb P\) for the integral group ring of its multiplicative group; the auxiliary addition is not imposed as a relation in this group ring.

Let \(\mathcal A\) be a cluster algebra over \(\mathbb P\), with coefficient ring \(\mathbb Z\mathbb P\). For a cluster \(\mathbf x=(x_1,\ldots,x_n)\), consisting of its mutable cluster variables, a cluster monomial is an element
\[
 x_1^{a_1}\cdots x_n^{a_n},\qquad (a_1,\ldots,a_n)\in\mathbb Z_{\geq 0}^{n}.
\]
Let \(\operatorname{CM}(\mathcal A)\) be the set of these elements as the cluster and exponents vary. Equal elements obtained from different clusters are counted only once. Coefficients belong to the scalar ring rather than being additional factors indexing distinct cluster monomials.

\textbf{Question.} For every coefficient semifield \(\mathbb P\) and every cluster algebra \(\mathcal A\) over \(\mathbb P\) covered by the source formulation, is \(\operatorname{CM}(\mathcal A)\) linearly independent over \(\mathbb Z\mathbb P\)? Explicitly, for every integer \(r\geq 1\), every choice of pairwise distinct \(m_1,\ldots,m_r\in\operatorname{CM}(\mathcal A)\), and every \(c_1,\ldots,c_r\in\mathbb Z\mathbb P\), must
\[
 \sum_{j=1}^{r}c_jm_j=0\quad\Longrightarrow\quad c_1=\cdots=c_r=0
\]
hold in \(\mathcal A\)?

\textbf{Unresolved scope qualifications.} The extracted statement does not specify the permitted exchange matrices or explicitly fix the cluster-algebra convention. The definition above uses the usual convention of finite clusters and mutable-variable cluster monomials, with coefficients treated as scalars. No skew-symmetry restriction is added to the question. The intended seed class and coefficient conventions require confirmation against the source.

The supplied dated curation labels the problem partially solved and attributes a linear-independence result for all skew-symmetric cluster algebras to the cited primary proof. It also states that the extracted formulation permits arbitrary coefficient semifields and that the cited result does not establish the full stated generality. These are supplied source assessments, not independently verified mathematical or literature conclusions.
\subsection{Short English statement}
The question asks whether every finite relation among distinct cluster monomials, with coefficients in \(\mathbb Z\mathbb P\), has all coefficients zero. The supplied formulation imposes no explicit skew-symmetry restriction. Its precise seed and coefficient conventions need review, and the supplied partially solved status is a dated curation assessment.
\subsection{Sources}
\begin{itemize}
\item Local UnsolvedMath registry, permanent ID \(30002455\), title \emph{Linear Independence of Cluster Monomials}, source problem code \texttt{catalogue entry}, in the supplied Oberwolfach Reports Open Problems collection. The supplied record gives proposed year 2013 and names no proposer.
\item Supplied source citation: \emph{Cluster Algebras and Related Topics (2014)}, Oberwolfach Reports, DOI \texttt{10.4171/owr/2013/58}, pp.~3379--3432. Original OWR record: \url{https://doi.org/10.4171/owr/2013/58}. The background describes these reports as records of mathematical workshops and their open-problem discussions.
\item Supplied literature source, labeled ``Cerulli Irelli Keller Labardini Fragoso Plamondon'': \url{https://arxiv.org/abs/1203.1307}. The supplied assessment attributes linear independence for all skew-symmetric cluster algebras to this proof, while stating that it does not establish the extracted question in its full generality with arbitrary coefficient semifields.
\item Literature assessment date: \texttt{2026-08-21}. This is supplied dated curation, not independently verified. The background expressly requests independent verification of time-sensitive status and cited claims.
\item The registry labels the statement \texttt{corrected\_verified} and reports that it was checked against the official Oberwolfach report. This verification claim is supplied metadata and has not been independently checked here. The supplied problem status is \texttt{partially\_solved}.
\end{itemize}
\end{document}
